\section{Introduction}
\label{sec:introduction}

Offline reinforcement learning (RL) learns from a fixed dataset and cannot use new interaction to correct extrapolation errors. Behavior-regularized actor--critic methods therefore face a coupled design problem: policy improvement should exploit the critic, yet the actions entering Bellman targets should remain in regions where that critic is reliable. In TD3+BC, the same actor is behavior-regularized, copied into the target network, and evaluated at deployment \citep{fujimoto2021minimalist}. Increasing the actor's value coefficient consequently changes both the target actions used to train the critic and the reach of the deployed policy.

We examine this coupling through \emph{budgeted actor refinement} (BAR). At every delayed policy update, BAR maintains $K$ persistent deterministic actors. The first actor performs the behavior-anchored TD3+BC update; later actors are anchored to the stop-gradient output of their predecessor. A Polyak copy of the first actor enters the TD target, whereas the final actor is evaluated. For a total nominal coefficient budget $T$, each actor uses $h=T/K$. Here ``budget'' means the sum of nominal proximal coefficients in the actor objectives. It is not a wall-clock, optimizer-call, or displacement budget: increasing $K$ also increases actor compute and parameters, and the realized movement depends on the critic, clipping, and optimization trajectory.

Our primary object is a \emph{stability phase diagram}, not a selected benchmark point. We operationally call the region stable when, on the tested environment--budget grid, the two-seed mean D4RL score is at least 20; ``high budget'' means $T\geq4$. This threshold is descriptive and is varied in the supplement. Across proximal configurations $K=1,2,3,4$, the high-budget mean increases from $25.97$ to $42.02$, $52.78$, and $63.61$, while collapsed cells fall from $35/63$ to $25/63$, $17/63$, and $9/63$. Dataset-cluster resampling supports the broad $K=4$ versus $K=1$ rescue, whereas the reduction in raw collapse risk from $K=3$ to $K=4$ has an interval that reaches zero. Thus the strongest evidence is a broad refinement-depth effect relative to one-step TD3+BC; the incremental $K=4$ versus $K=3$ failure-rate claim is more tentative.

The geometric analysis is deliberately narrower than the empirical claim. For deterministic policies, squared action anchoring is a state-fiber Wasserstein--2 cost. This yields an ideal proximal operator, a conditional learned-critic descent inequality, and a first-hop target-perturbation bound. The live implementation, however, uses independently initialized persistent actors and only one Adam update per hop. It need not realize a near-identity proximal map, and the analysis does not certify its coupled training dynamics. We therefore use geometry as a local organizing principle, not as a proof of the observed large-budget stability.

Our contributions are:
\begin{itemize}
\item We introduce a persistent TD3+BC actor chain that separates the target-facing first branch from the deployed final branch while distributing a fixed nominal actor coefficient across refinement depth.
\item We distinguish the implemented one-update actor chain from its ideal deterministic Wasserstein proximal interpretation, and provide conditional finite-step and target-perturbation statements with their precise scope.
\item We report a complete fixed-budget phase diagram through $K=4$, dataset-cluster uncertainty, a projected-linearized control, an independent one-hop seed replication, and compact re-audited diagnostics of target movement and critic failure.
\end{itemize}
